% !TEX root = ../main.tex

\ifprmresultsready

%----------------------------------------------------------------------
\subsection{Physical reading of the error structure}
%----------------------------------------------------------------------

The transfer hierarchy has a structural reading.  Recombining familiar
constituent representations is nearly free; the expensive step is a new
constituent, and the preregistered novelty analysis attributes the
host--impurity asymmetry: the sample-matched error gap follows the
structural novelty of the held-out host --- its composition, cell,
thickness, and coordination environment relative to the training hosts ---
and survives controlling for the mechanical training-support asymmetry
between the two protocols.  A new host changes the structural motif the
defect embeds into, whereas a new impurity largely re-poses a chemistry
seen in other combinations.  Development-era case studies of five named
held-out hosts (Supplemental Material~\cite{supplemental}) trace the same
gradient concretely: hosts with close structural relatives in training
transfer at no measurable cost, while the one host with no structural
relative in the database sits at the failure extreme.

These grouped tests do not represent element- or formula-disjoint
pretraining.  The source task contains reduced-formula matches for 26 IMP2D
host labels and every IMP2D impurity element, so host- and impurity-held-out
mean absence from formation-energy supervision, not absence from all earlier
representation learning.  All conclusions are internal to the retained 44
hosts and 65 impurities --- a set that contains no nitride host --- and to
the IMP2D reference conventions.

The class asymmetry is likewise not one phenomenon.  A snugly caged
impurity sits in a steep, sterically pinned energy well whose depth
follows the local geometry the model sees directly; a loosely held
impurity in a large, uniform cage samples a flatter landscape with
near-degenerate placements, where the energy depends on subtler electronic
contributions --- which is why interstitial error concentrates in roomy,
loosely constrained cages while adsorbate error tracks chemical mismatch.
These are controlled associations, not causal mechanisms; the sign pattern
is what the frozen analysis established.  The measured descriptors
attenuate the class gap by roughly a quarter at the point estimate, but
the attenuation interval includes zero, so the penalty is reported as not
yet explained: the remainder must live in physics the frozen descriptors
do not span, such as the anisotropy of the relaxation pattern, which the
released database does not record.  The impurity-identity audit also makes
the retained population selective (316 interstitial and 33 adsorbate rows
removed), so the asymmetry applies to the retained population rather than
to every released IMP2D row.

The readout result follows the same defect-local logic.  The formation
energy of Eq.~\eqref{eq:formation_energy} is a defect excess energy --- a
difference of extensive totals that does not grow with host atom count ---
so atomwise additive pooling imposes the wrong scaling class, and pooling
choice is known to matter for intensive and localized
targets~\cite{buterez2023pooling}.  The extensivity diagnostic supports
this mechanism at the residual level, with a positive additive
residual-size slope whose host-clustered interval nevertheless spans zero,
so the mechanism is suggested rather than established; development-era
attribution analyses (Supplemental Material~\cite{supplemental}) show the
same property from inside the model, with occlusion attributions
concentrated at the impurity atom and decaying with distance from it.
These are complete-recipe contrasts under one protocol, not a universal
backbone ranking.

%----------------------------------------------------------------------
\subsection{Scope of the reranking claim}
%----------------------------------------------------------------------

For a fixed host--impurity pair, the pristine-host energy and impurity
chemical potential in Eq.~\eqref{eq:formation_energy} cancel when sites or
incorporation classes are ranked, which is why useful relative decisions
coexist with a larger absolute error.  The decisions are conditional on
the DFT-relaxed candidates released by IMP2D: each geometry and its
reference energy are outputs of the same completed workflow, so no
relaxation cost is bypassed, and the results concern retrospective energy
estimation and reranking --- not structure generation, relaxation,
thermodynamic populations, kinetics, or synthesizability.

%----------------------------------------------------------------------
\subsection{Correctness, uncertainty, and remaining scope}
%----------------------------------------------------------------------

Every number in this paper comes from a repaired representation whose
permutation invariance and exact minimum-image distances passed an
independent acceptance audit; the source-task pretraining was rebuilt under
the same corrected construction.  The diagnostic of the superseded
implementation is a result in its own right [Fig.~\ref{fig:permutation}]:
sixteen atom-relabeling permutations moved single predictions by up to
$8.8$\eV{} and flipped incorporation-class decisions, while the pooled MAE
moved by less than $10^{-3}$\eV{}.  Aggregate benchmarks are therefore
nearly blind to representation defects that dominate individual predictions
--- a caution we believe generalizes beyond this system.

\begin{figure}[t]
\centering
\includegraphics[width=\columnwidth]{fig_permutation_repair.pdf}
\caption{Atom relabeling as a representation diagnostic.
  (a) Cumulative distribution of the per-structure prediction spread of the
  superseded (legacy) implementation across sixteen atom permutations, by
  incorporation class ($p_{99}=0.46$\eV{}, maximum $8.8$\eV{}); under the
  repaired representation every structure reproduces its prediction to
  within the archived round-trip bound ($\leq 1.7\times10^{-5}$\eV{},
  vertical line).
  (b) The same axis at aggregate level: the legacy pooled MAE changes by at
  most $8\times10^{-4}$\eV{} across permutations --- four orders of
  magnitude below the largest single-structure spread --- while individual
  relabelings flip 9--19 incorporation-class decisions and 76--102 site
  selections.  Aggregate benchmarks are nearly blind to this defect.}
\label{fig:permutation}
\end{figure}

The evidence remains limited to neutral adsorbates and interstitials, the
source database's relaxed structures and reference energies, and a model
that infers the reference-energy difference from the impurity-containing
graph alone.  Vacancies, substitutions, charged defects, transition levels,
migration barriers, finite-temperature free energies, and competing phases
are outside scope.  Duplicate control removes numerical equivalents
detected by the two declared structural fingerprints but cannot prove that
every physical redundancy has been found.  SchNet and the descriptor
families are aligned comparators under the frozen protocol, not an
exhaustive survey; a development-era survey (Supplemental
Material~\cite{supplemental}) additionally evaluated trained equivariant
models (MACE, ViSNet), ALIGNN, and the MACE-MP-0 foundation model as a
zero-shot label surrogate --- the latter essentially uncorrelated with the
IMP2D formation-energy convention --- and none approached the
defect-centered design at matched scale, consistent with reports that
universal interatomic potentials degrade on interstitial-like environments
far from their bulk training distributions.  In absolute terms, the
repaired random-split and pair-held-out errors sit in the
few-tenths-of-an-eV band reported for specialized defect formation-energy
models on bulk and low-dimensional
benchmarks~\cite{witman2023,kiyohara2025,kazeev2023}; the grouped regimes
report the additional cost of chemical novelty that pooled numbers do not
expose.  Internal ensemble uncertainty (Supplemental
Material~\cite{supplemental}) supports marginal calibration and risk
ranking within the retained population, not host-conditional or chemically
shifted coverage.

The single-database scope is itself an evidence-based protocol choice.
Development-era experiments (Supplemental Material~\cite{supplemental})
found zero-shot cross-database evaluation on JARVIS vacancy data failing
at the 2--3\,eV level with negative $R^2$; few-shot fine-tuning lost to
training from scratch; and multi-source joint training bought small
in-distribution gains at the price of degraded held-out-host transfer.
The barrier is one of reference conventions and defect-class mismatch.
The most direct next test is therefore a target-matched external set:
structures, chemical potentials, convergence settings, and target
definition matching the IMP2D convention, generated independently of their
reference energies.  A development-era pilot (Supplemental
Material~\cite{supplemental}) has already exercised that test end to end
with independent plane-wave calculations on model-ranked candidates; it
located the barrier in per-element reference conventions and confirmed
that model ranking concentrates first-principles effort well above the
pool base rate.  The remaining step is convention matching, not new
machinery.

\else

\noindent\textit{The numerical discussion is unavailable because the audited
result bundle is not present.}

\fi
